% Related work. Part of the team paper; included by main.tex.
% Every citation here was written from memory: check each claim against the paper before submission (%CHECK).
\section{Related work}
\label{sec:related}

Our work combines three lines of research: controlled studies of in-context learning, transformers for sequential decision
making, and imitation from a few demonstrations. Each subsection ends with how our setup differs.

\subsection{In-context learning as a controlled phenomenon}
In-context learning was first observed at scale in GPT-3~\cite{brown2020language}, which performs new tasks from a handful of
examples in its prompt. Because natural-language tasks have no known ground truth, later work moved to synthetic settings where
the task family is fixed by design. Garg et al.~\cite{garg2022can} trained transformers from scratch on prompts of input--output
pairs from linear functions, sparse linear functions, decision trees and two-layer networks. In-context predictions matched the
optimal estimator (least squares for linear functions) and were robust to some distribution shifts between training and test
prompts. Several papers then asked \emph{which} algorithm the transformer implements. Aky{\"u}rek et al.~\cite{akyurek2023learning}
showed that transformers trained on linear regression behave like gradient descent, ridge regression or exact least squares
depending on depth and width. Von Oswald et al.~\cite{vonoswald2023transformers} constructed linear self-attention layers that
perform a gradient-descent step and found that trained models match them.

Other work studies \emph{when} in-context learning appears. Chan et al.~\cite{chan2022data} found that it depends on properties of
the training data, such as bursty sequences and many rare classes, and competes with storing knowledge in the weights. Raventós et
al.~\cite{raventos2023pretraining} showed for linear regression that below a threshold of task diversity, transformers behave like a
Bayesian estimator with the pretraining tasks as prior and fail on new tasks. Above it, they generalize to new tasks like ridge
regression. Xie et al.~\cite{xie2022explanation} interpret in-context learning as implicit Bayesian inference of a latent concept
shared by the prompt. Mechanistic work identified induction heads, attention circuits that copy patterns from earlier in the
context, as one mechanism behind in-context learning in language models~\cite{olsson2022context}.

\textbf{Relation to our work.} We keep the Garg et al.\ recipe: a GPT-2 trained from scratch, a known task family, prompts of
examples followed by a query, and evaluation on held-out tasks. The ``function'' is now a goal-conditioned policy with the goal
hidden, and the ``examples'' are trajectories. We add closed-loop evaluation, where predictions change future inputs, which these
studies do not have. The Bayesian view fits our setting well. With no demonstrations, the best prediction averages over goals,
our no-context baseline. With demonstrations, the model can concentrate on the demonstrated goal. The task-diversity question
of~\cite{raventos2023pretraining} is a natural next experiment for us.

\subsection{Transformers for sequential decision making}
Decision Transformer~\cite{chen2021decision} and Trajectory Transformer~\cite{janner2021offline} cast offline reinforcement learning
as sequence modeling over states, actions and returns. Gato~\cite{reed2022generalist} trained one transformer on many control and
language tasks and could be prompted with demonstrations of a task. Prompt-DT~\cite{xu2022prompting} extends Decision Transformer to
few-shot generalization: a short trajectory segment from a new task, placed in the prompt, specifies the task to solve.

A second group of papers trains transformers to learn \emph{reinforcement learning} in context. Algorithm
Distillation~\cite{laskin2023incontext} trains on the learning histories of an RL algorithm, so the model improves its policy within
its context across episodes. The Decision-Pretrained Transformer~\cite{lee2023supervised} is trained to predict the optimal action
for a query state given an in-context dataset of interactions, and it explores and exploits in new tasks. Raparthy et
al.~\cite{raparthy2023generalization} trained transformers to solve new sequential decision tasks from a few demonstrations and
studied which data properties make this possible.

Closest to our motivation, Smith et al.~\cite{smith2026transformers} study in-context \emph{low-level} control. A GPT-2-style
transformer stabilizes cartpole and acrobot systems whose physical parameters are hidden, inferring them from a short unforced
trajectory in its context while imitating an analytic expert controller. Success generally improves with context length, and some
out-of-distribution systems are stabilized. Training on noise-free expert trajectories was a key failure: in-distribution cartpole
success fell to 5\%, from 86\% when states were perturbed during expert control.

\textbf{Relation to our work.} Prompt-DT, Gato and Raparthy et al.\ show that trajectory prompts can specify a task, but in large
multi-task settings where context use is hard to isolate. Smith et al.\ show in-context control with hidden dynamics. Our context
instead holds demonstrations of the task, the hidden variable is the goal rather than physical parameters, the task involves
contact, and wrong-goal controls test whether the model uses its context. Algorithm Distillation and DPT learn from rewards. Our setting uses
no rewards and no return conditioning: the model imitates fixed teachers, and the only task variable is a hidden goal. That makes
the controls exact, since a wrong-goal prompt changes only the goal and keeps the query states identical. It also lets us measure
directly how much of the teacher's success the in-context learner recovers.

\subsection{Imitation from a few demonstrations}
One-shot imitation learning~\cite{duan2017oneshot} is the closest prior setup. A network conditioned on one demonstration of a
block-stacking task, through attention over the demonstration and the current state, performs that task from new initial states,
including stacking orders unseen in training. Other approaches to imitation from a single demonstration
include meta-learning a policy that adapts by gradient steps~\cite{finn2017oneshot}, learning a task embedding from
demonstrations~\cite{james2018taskembedded}, and transformer attention between a demonstration video and the robot's
observations~\cite{dasari2020transformers}. Recent work brings in-context imitation to real robots: Keypoint Action
Tokens~\cite{dipalo2024keypoint} prompt an off-the-shelf language model with demonstrations encoded as tokens, and In-Context Robot
Transformer~\cite{fu2024incontext} trains a transformer by next-token prediction on prompts of teleoperated trajectories.

Meta-reinforcement learning infers a task from experience rather than demonstrations. RL$^2$~\cite{duan2016rl2} and Wang et
al.~\cite{wang2016learning} train recurrent policies whose hidden state adapts to a new task within a few episodes.
MAML~\cite{finn2017model} learns an initialization that adapts with a few gradient steps. PEARL~\cite{rakelly2019efficient} encodes
collected experience into a latent task variable that conditions the policy, the same role the hidden goal plays for us.

\textbf{Relation to our work.} Our model is a demonstration-conditioned policy, like one-shot imitation, but it uses a
general-purpose GPT-2 with no architecture specific to the task. We also make three measurements these systems usually do not
report together: performance against the goal-seeing teacher, wrong-goal and no-context controls on identical starts, and a
held-out region of task space.

\subsection{Imitation learning and compounding errors}
Behavior cloning~\cite{pomerleau1989alvinn} trains a policy by supervised learning on expert states. Its errors compound: the
learner drifts into states the expert never visited, and its error can grow quadratically with the horizon~\cite{ross2010efficient}.
DAgger~\cite{ross2011dagger} fixes this by labeling the learner's own states with the expert. DART~\cite{laskey2017dart} instead adds
noise to the expert's actions during data collection, so the data covers nearby states while keeping the expert's labels. Policy
distillation~\cite{rusu2016policy} compresses a trained RL policy into a smaller network by supervised learning. Push-T itself comes
from the implicit behavior cloning work of Florence et al.~\cite{florence2021implicit} and is a standard benchmark for imitation, for
example in Diffusion Policy~\cite{chi2023diffusion}.

\textbf{Relation to our work.} Our BC teacher is a policy distillation of the PPO teacher with DART-style noise injection. The
transformer is itself behavior cloning, conditioned on the context, and it shows the same compounding-error behavior: an offline
error 4$\times$ that of the BC teacher gave 0\% closed-loop success. Its training queries come only from clean teacher rollouts. Noise
injection on query states is the corresponding fix, and it is our next experiment.
